\documentclass[11pt]{article}

\usepackage[margin=1.1in]{geometry}
\usepackage{amsmath,amssymb}
\usepackage{booktabs}
\usepackage{array}
\usepackage{microtype}
\usepackage{listings}
\usepackage[hidelinks]{hyperref}
\usepackage{url}

\lstdefinestyle{pseudocode}{
  language=Python,
  basicstyle=\ttfamily\footnotesize,
  keywordstyle=\bfseries,
  commentstyle=\normalfont\itshape\footnotesize,
  columns=fullflexible,
  keepspaces=true,
  showstringspaces=false,
  frame=lines,
  aboveskip=0.7\baselineskip,
  belowskip=0.7\baselineskip
}

% --- notation -------------------------------------------------------
\newcommand{\chop}[1]{\mathrm{chop}_{#1}}   % truncate to #1 significant bits
\newcommand{\rn}[1]{\mathrm{RN}_{#1}}       % round-to-nearest-even at #1 bits
\newcommand{\ulp}{\mathrm{ulp}}
\newcommand{\instr}[1]{\texttt{#1}}
\newcommand{\hex}[1]{\texttt{#1}}

\title{Thirty Years of the Pentium Kernel:\\
A Bit-Exact Behavioral Reconstruction of the x87\\
Transcendental Instructions in Intel Skylake}

\author{}

\date{September 2026}

\begin{document}
\maketitle

\begin{abstract}
We reconstruct general, bit-exact numerical programs for standalone x87
\instr{FSIN}/\instr{FCOS} and paired \instr{FSINCOS} on an Intel Skylake Xeon.
The programs combine
exact integer reduction using the architectural 66-bit approximation to
$\pi$, the Pentium coefficient ROM and table cells, fixed internal rounding
operations, and an explicit tiny-result predecessor rule. If $T_p$ truncates
to $p$ significant bits, the ordinary multiply is
$M(x,y)=T_{67}(T_{67}(x)T_{64}(y))$; coefficient additions round to nearest-even
at 64 bits. Sine and cosine have distinct terminal schedules, and one table
multiply rounds its exact port product directly to 64 bits.
Paired \instr{FSINCOS} instead uses a Horner polynomial, truncating every
coefficient-stage product to 67 bits before its RN64 addition.
No captured-operand ledger, fitted exception selector, or rounding-history
correction participates in the promoted programs.
Standalone results match all 81 previously failing external frontier rows and all
3{,}379{,}017 retained output appearances in the promotion regression, including
3{,}378{,}987 applicable \texttt{C1} checks. Independent rational/integer
implementations, domain bounds, compiler-variant checks and previously frozen
one-shot hardware challenges support the reconstruction.
The paired program matches 23{,}838{,}534 retained lane results and all
13{,}800 fresh adversarial RC/PC tuples; its predecessor fails 750 of those
tuples. The default C implementation includes all three instructions.
These results establish validated behavioral algorithms, not a recovered
physical netlist or an exhaustive enumeration of x87 inputs.
\end{abstract}

\section{Introduction}

Every Intel processor since the 8087 has offered hardware
transcendental instructions.  Their architectural documentation is
thin: the Software Developer's Manual specifies operand ranges, flag
behavior, and---since a 2014 controversy over \instr{FSIN} accuracy
claims~\cite{dawson2014,intel-sdm}---the fact that range reduction
uses a 66-bit approximation to $\pi$.  The architectural description
does not specify the kernels, tables, working precisions, rounding of
each internal operation, or exact behavior at every boundary.
This gap has practical consequences.  Emulators and binary-translation
systems (QEMU, Bochs, Valgrind, hypervisor instruction decoders)
approximate these instructions with host arithmetic or software
libraries, so guest-visible results differ from silicon; record-replay
debugging, live-migration between vendors, deterministic sandboxing,
and archival execution of legacy x87 code all inherit the discrepancy.
The algorithms in shipping silicon can be studied through public
implementation material and measured processor behavior.

This paper presents the validated standalone and paired numerical programs and their
implementation, retaining the shared-kernel and ROM-lineage results as
context. The target is one Skylake Xeon reference, not universal parity across
Intel and AMD processors. Earlier R96-stage models used an incomplete support
selector and, in some comparisons, a captured-output overlay. They are not
the algorithm now promoted.

\subsection*{Contributions}

\begin{enumerate}
\item A complete numerical specification of the standalone polynomial,
table, reduction and tiny paths, with explicit working precisions and
rounding points (Section~\ref{sec:algorithm}).
\item A default C implementation of that fixed program, resolving all
81 recorded incumbent-frontier rows without an operand ledger, and a
retained regression separating output, \texttt{C1}, prospective hardware
evidence and software equivalence (Section~\ref{sec:validation}).
\item Domain arguments explaining the asymmetric multiply representation
and exact table-cell selection, plus negative controls distinguishing
numerically different evaluation schedules.
\item An independently reconstructed paired direct-polynomial Horner program, validated by
retained regression and a frozen all-RC, PC24/53/64 adversarial challenge
(Sections~\ref{sec:paired} and~\ref{sec:paired-validation}).
\item Behavioral evidence for Pentium ROM lineage and corrections to the
published ROM decode (Section~\ref{sec:lineage}), with physical-routing
inferences kept separate from numerical results.
\end{enumerate}

\section{Background and related work}
\label{sec:background}

\paragraph{The x87 model.}  The x87 unit computes in 80-bit double
extended precision: a 64-bit explicit-integer-bit significand, 15-bit
exponent, sign.  The transcendental instructions consume the top of
the register stack and are sensitive to the rounding-control (RC)
field of the control word; \instr{FSIN}/\instr{FCOS}/\instr{FSINCOS}
accept $|x| < 2^{63}$ and otherwise set the \texttt{C2} flag leaving
the operand unchanged; \instr{FPTAN} additionally pushes a constant
$1.0$; \instr{F2XM1} computes $2^x-1$ on $|x| \le 1$.

\paragraph{History.}  The 8087/80287 used CORDIC-style iteration; the
80387 added \instr{FSINCOS}.  The Pentium (1993) moved to a
table-assisted polynomial design whose 304-entry constant ROM was
physically decoded from die photographs by Shirriff in
2025~\cite{shirriff2025}.  For Itanium, Intel published and formally
verified a table-free double-extended algorithm intended also to serve
IA-32 compatibility~\cite{harrison2000,harrison1999}; its assembly
source survives in the glibc ia64 port (removed upstream in
glibc~2.39), from which we transcribed a frozen bit-exact reference
model.  We compare these published lineages with the numerical behavior
of the selected Skylake reference.  AMD's K5
implementation was described by its designers~\cite{lynch1995}, and
folklore has long held that AMD's argument reduction is more accurate
than Intel's.  The 2014 accuracy controversy~\cite{dawson2014} led
Intel to document the 66-bit-$\pi$ reduction property, but not the
algorithm.

\paragraph{Related reverse engineering.}  The Pentium \instr{FDIV}
defect of 1994~\cite{nicely1994,coe1995} established that arithmetic
microstructure (there, an SRT quotient-digit table) can be recovered
from behavior; Shirriff's die-level work recovered the Pentium FPU's
constants and layout~\cite{shirriff2025}.  This work specifies and tests
the standalone and paired numerical paths on the selected Skylake reference.
It makes no claim to be the first implementation or study of these instructions.

\section{Methodology}
\label{sec:method}

The numerical target is the Skylake Xeon reference used for the frozen
standalone and paired campaigns. Ground truth consists of raw instruction outputs and
status words captured on real x86 hardware, not host-library trigonometric
values. The four rounding modes are RN, RD, RU and RZ. Where precision-control
independence is tested, PC24, PC53 and PC64 are separate recorded tuples.

A hypothesis is specified before fresh hardware observation. Proposed inputs
are checked against retained capture history, predictions and manifests are
frozen, and each eligible tuple is executed once. Existing observations are
reused for regression; they are never called fresh merely because a new
implementation scores them. Tests include model-disagreement operands and
common-prediction controls near rounding and dispatch boundaries.

Independent verification uses separately implemented rational or integer
arithmetic, reduction, quantization and raw-output parsing. Compiler comparisons
check faithful realization of the same specification; they are not additional
hardware observations. Counts of retained bank appearances can overlap and
are not summed into a unique-input or fresh-capture claim.

Here \emph{bit-exact} describes agreement of the behavioral program with the
specified reference observations. Generality is supported by fixed operators,
domain reasoning, adversarial transfer and implementation checks, not by an
input-indexed exception table. Numerical equivalence does not identify hidden
physical operand-port routing. Neither exhaustive enumeration nor netlist
recovery is required to state a validated behavioral algorithm.

\section{The algorithm}
\label{sec:algorithm}

For a code-first walkthrough, Appendix~\ref{app:pseudocode} gives
programmer's pseudocode for reduction, dispatch, both polynomial schedules,
the shared table kernel and final rounding. Its helper functions correspond
directly to the equations in this section.

Throughout, $\chop{k}(\cdot)$ truncates a product to $k$ significant
bits (magnitude chop); $\rn{k}(\cdot)$ rounds to nearest-even at $k$
bits.  The datapath evaluates significands as integers; all laws below
are integer/bit-level statements about those significands.

\subsection{Encoding classes and specials \textnormal{[EXACT]}}
\label{sec:specials}

Unnormals (nonzero exponent, explicit integer bit 0), pseudo-NaNs and
pseudo-infinities take the masked invalid-operation response: both
results are the real indefinite \hex{ffff:c000000000000000}, in every
rounding mode.  Pseudo-denormals (exponent 0, integer bit 1) are valid
and compute normally.  NaNs are quieted and passed through;
infinities produce the indefinite; zeros are exact ($\sin\pm0 =
\pm0$, $\cos\to 1$).  $|x| \ge 2^{63}$ sets \texttt{C2} with the
operand unchanged, exactly at the encoding boundary. (Historical i7
invalid-encoding validation: 1{,}408/1{,}408.)

\subsection{Exact integer range reduction}
\label{sec:reduction}

Write a finite normalized magnitude as $|x|=s\,2^{e-63}$, where
$2^{63}\le s<2^{64}$. Inputs below the direct-path threshold
\hex{3ffe:c90fdaa22168c234} use the unreduced argument.
For the reduction domain $-1\le e\le62$, define
\[
 M_{66}=\hex{3243f6a8885a308d3}_{16},\qquad
 A=s\,2^{e+2},\qquad
 q=\left\lfloor\frac{A}{M_{66}}+\frac12\right\rfloor,\qquad
 D=A-qM_{66}.
\]
The signed residual is $(-1)^{\sigma_x}D\,2^{-65}$.
Thus the reduction constant is $(\pi/2)_{66}=M_{66}2^{-65}$.
Since $M_{66}$ is odd, the quotient has no half-integer tie.
The implementation's two-part residual reconstructs this exact dyadic
difference before kernel dispatch; no quotient or carry history is appended
to the numerical value.

For external phase $\phi=0$ (\instr{FSIN}) or $\phi=1$ (\instr{FCOS}), let
$n=(-1)^{\sigma_x}q+\phi$; on the direct path, use $q=0$.
Let $r$ be the residual magnitude and $\sigma_r$ its sign.
The selected internal function is cosine if $n\bmod2=1$ and sine otherwise.
The result sign is bit 1 of $n$, XOR $\sigma_r$ only on the sine branch.

Integer interval and modular arguments cover the full normalized reduction
domain. In particular $D\ne0$: coprimality would otherwise require
$M_{66}\mid s$, impossible for a nonzero 64-bit $s$. The small nonzero
residuals near multiples of the reduction constant therefore have an explicit
tiny rule rather than a guessed zero-residual case.

\subsection{Path selection and arithmetic operators}

For $r>0$, let $E=\lfloor\log_2 r\rfloor$.
The tiny path has $E<-32$, the polynomial path $-32\le E\le-3$,
and the table path $E\ge-2$. All comparisons are exact binary comparisons.
Write $T_p=\chop{p}$ for signed magnitude truncation. For the standalone
polynomial and shared table programs, define
\[
 \begin{aligned}
 M(x,y)&=T_{67}\bigl(T_{67}(x)\,T_{64}(y)\bigr),\\
 A_{64}(x,y)&=\rn{64}(x+y),\\
 M_R(x,y)&=\rn{64}\bigl(T_{67}(x)\,T_{64}(y)\bigr).
 \end{aligned}
\]
The product inside $M_R$ is exact before its one RN64 rounding:
$M_R(x,y)$ is \emph{not} $\rn{64}(M(x,y))$.
Final $\mathrm{RC}_{64}$ applies the requested architectural rounding mode
to the signed prevalue.

\subsection{The fixed standalone polynomial program}
\label{sec:polynomial}

Let $K_1,\ldots,K_6$ be the selected native sine or cosine coefficients.
Both families use the same graph:
\[
 \begin{aligned}
 S&=M(r,r),& F&=M(S,S),\\
 N&=A_{64}\bigl(K_1,M(F,A_{64}(K_3,M(F,K_5)))\bigr),\\
 P&=A_{64}\bigl(K_2,M(F,A_{64}(K_4,M(F,K_6)))\bigr),\\
 L&=M(S,N),& R&=M(F,P).
 \end{aligned}
\]
With result sign $\sigma$ determined above, the terminals are
\[
 \begin{aligned}
 y_{\sin}&=\mathrm{RC}_{64}\left((-1)^\sigma
       \left[r+M(r,A_{64}(L,R))\right]\right),\\
 y_{\cos}&=\mathrm{RC}_{64}\left((-1)^\sigma
       \left[1+T_{67}(L+R)\right]\right).
 \end{aligned}
\]
The terminals are deliberately different; there is no learned condition
choosing which schedule applies to a known error operand.

Every direct polynomial residual has at most 64 significant bits and every
reduced polynomial residual at most 63. Multiplications produce at most 67
bits and Horner additions at most 64. With the native coefficient widths,
the only nonidentity input-port truncation in this graph is the fourth-power
producer's $T_{64}(S)$. This explains the numerical X67/Y64 representation
without claiming that a physical port has been observed.
The shared polynomial program replaces the old terminal-selector composition
in the promoted standalone route; it is not a new correction layered onto it.

\subsection{The fixed table program}
\label{sec:table}

For $r\ge1/4$, select the native center $b/64$ using
\[
 b=\begin{cases}
 18+4\lfloor16(r-1/4)\rfloor,&r<1/2,\\
 36+8\min(2,\lfloor8(r-1/2)\rfloor),&r\ge1/2.
 \end{cases}
\]
Only $b\in\{18,22,26,30,36,44,52\}$ is reachable.
Put $a=r-b/64$ and $S=M(a,a)$.
For each native four-coefficient family define
\[
 H(K)=A_{64}\bigl(K_1,M(S,A_{64}(K_2,
             M(S,A_{64}(K_3,M(S,K_4)))))\bigr).
\]
Let $p=H(K^{\sin})$, $q=H(K^{\cos})$, and let $t_s,t_c$ be the
native sine/cosine ROM entries at the selected center. Then
\[
 \begin{aligned}
 v&=A_{64}(a,M(M(S,p),a)),& w&=M_R(S,q),\\
 z_s&=T_{67}\bigl(M(t_c,v)+M(t_s,w)\bigr),\\
 z_c&=T_{67}\bigl(-M(t_s,v)+M(t_c,w)\bigr),\\
 y_{\sin}&=\mathrm{RC}_{64}\bigl((-1)^\sigma(t_s+z_s)\bigr),\\
 y_{\cos}&=\mathrm{RC}_{64}\bigl((-1)^\sigma(t_c+z_c)\bigr).
 \end{aligned}
\]
The validated sine $K_4$ adjustment, subtracting $2^{-25}$ from
the decoded coefficient, is retained unchanged; no coefficient was fitted
to the newly resolved frontier.

Center subtraction gives at most 61 significant offset bits for reduced
inputs and at most 60 for direct inputs. All table input-port cuts are
identities under the displayed operand ordering. Thus this is a numerical
representation of the table program, not separate evidence for physical
table-port asymmetry. Exact center offsets annihilate the correction
products and expose the ROM word and final rounding directly. The $b=52$
center itself is outside the reachable residual domain; row $b=60$ is
unreachable.

The bit-slice classifier must not pass through a binary64 intermediate.
Such a model misclassifies a narrow sliver below interior boundaries;
retained blind boundary tests selected exact classification on all
7{,}936 operands, against 1{,}954 errors in that predecessor.

\subsection{Tiny values, final rounding and \texttt{C1}}
\label{sec:tiny}

For a non-bypass tiny residual, the leading magnitude is $r$ on the sine
branch and 1 on the cosine branch. RN and directed rounding away from zero
return that leading magnitude, with \texttt{C1}=1. Directed rounding toward
zero returns its 64-bit predecessor, with \texttt{C1}=0. The predecessor of
an exact power of two lies in the preceding binade and is handled explicitly.

The external-input bypass is distinct: a direct input with exponent below
$-68$ returns the input value for sine and $+1$ for cosine, with
\texttt{C1}=0, in all four modes. True denormals and pseudo-denormals take
this bypass; encoding-class information is kept separate from numerical
normalization.

For polynomial and table paths, \texttt{C1} records whether the final rounded
\emph{magnitude} exceeds the exact magnitude of the displayed terminal
prevalue. This is not a test for the entire instruction's precision exception.
In particular, exact final rounding does not imply that all earlier internal
operations were exact.

\subsection{Promoted entry points}
\label{sec:terminal}

The default standalone entry points invoke the fixed program above; the
paired entry invokes the distinct program below.
The historical R84 captured-operand overlay is disabled and unreachable
from these entry points. R96's empirical support selector and the later
experimental carry/history corrections are bypassed, not reinterpreted as
confirmed circuitry. Only the native ROM/coefficient tables and structural
encoding, reduction, quadrant, cell and tiny-path branches remain.

Paired \instr{FSINCOS} is not two independently validated standalone calls.
Retained Skylake data show standalone/paired disagreement on 987 of
240{,}000 dense-set operands. This establishes different numerical behavior,
not randomness or impossibility of an operand-dependent implementation.
The paired direct-polynomial program is independently reconstructed and
validated. Its difference from standalone does not require a statistical
approximation.

\subsection{The fixed paired polynomial program}
\label{sec:paired}

For $2^{-32}\le r<1/4$, let $K_i^s,K_i^c$ denote the same native
six-coefficient sine and cosine families. Compute the square once:
\[
 S=T_{67}(r^2),\qquad p_6=K_6^s,\qquad q_6=K_6^c.
\]
For $i=5,4,3,2,1$, evaluate
\[
 p_i=\rn{64}\bigl(T_{67}(p_{i+1}S)+K_i^s\bigr),\qquad
 q_i=\rn{64}\bigl(T_{67}(q_{i+1}S)+K_i^c\bigr).
\]
Every displayed product is materialized at 67 bits; its sum with the native
coefficient is exact before RN64 rounding. The schedule is uniform across
both arms and all five coefficient additions.
The two positive-residual prevalues are
\[
 u_s=r+T_{67}\bigl(\rn{64}(p_1S)\,r\bigr),\qquad
 u_c=1+T_{67}(q_1S).
\]
Apply the exact quadrant/residual-sign mapping separately to the two lanes,
then architectural $\mathrm{RC}_{64}$. The final external cosine lane supplies
\texttt{C1}, even when quadrant mapping selects internal $u_s$ for that lane.
The paired path shares the exact reduction, table, tiny and encoding rules;
it does not evaluate the standalone polynomial twice.

The cuts are observable arithmetic, not alternative notation for fused
multiply-adds. For RN input \texttt{3ffc:e79000000c3e46e7}, the last-product-only
schedule gives a sine ending in \texttt{d857}; hardware and the all-product
schedule give \texttt{d858}. Both share the prefix \texttt{3ffc:e5980e1fae54}.
The first stored
sine difference occurs at the $K_2^s$ addition. Exact replay places the two
prevalues respectively $5/1024$ ulp below and $5/512$ ulp above the final RN
midpoint. The cosine lane and directed-rounding outputs agree. Thus this
operand rejects the last-product-only schedule without an operand selector.

Exact rational interval enclosures cover all 30 polynomial binades. Across
840 operation-bound checks, both signed addends, products and final C1
re-encodings need at most 139 accumulator bits and alignment shifts at most
132, within the signed 256-bit implementation. These are source-backed
transcription bounds under the established reducer and quantizer contracts,
not an automatically verified C translation or a silicon theorem.

The promoted all-product representation specifies validated numerical
behavior, not uniquely recovered physical multiplier ports or a netlist.

\subsection{\instr{FPTAN} \textnormal{[EXACT]}}

\instr{FPTAN} computes a sine/cosine pair with its own terminal profile and
one final divide, then pushes $1.0$. Recovered operation profile: lookup
constants are 67-bit; the cosine tail is read at $\rn{64}$; the sine
analysis carrier is 69-bit, read through $\rn{64}$ before the
multiplier; ordinary state-producer multiplies are $\chop{67}$.
NaN and infinity operands push a \emph{second copy of the special
result} rather than $1.0$.  The table-lane classifier carries exactly
the Section~\ref{sec:table} structure.  Validation: zero mismatches
on a $2^{30}$-input exhaustive traversal of the binary64-embedded
space ($3.2\times10^{9}$ results, both \texttt{C1} flag states) plus
the structured probe sets.

\subsection{\instr{F2XM1} \textnormal{[EXACT]}}

$2^x - 1$ on $|x| \le 1$ is a separate polynomial engine with the
same datapath dialect: input exponent below $-68$ takes a final
$\ln 2 \cdot x$ multiply; exponents $-68$ up to $1/4$ take an
eleven-coefficient polynomial in $z = \ln 2 \cdot x$; $1/4 \le |x| <
1$ reduces around a midpoint $b$ and evaluates a six-coefficient
polynomial in $z = \ln 2\,(x - b)$, forming $(2^b - 1) + 2^b\,
\mathrm{expm1}(z)$.  The midpoint numerators over 128 are the odd
values $33..63$ for $|x| \in [1/4, 1/2)$ and $66, 70, \ldots, 126$
for $|x| \in [1/2, 1)$; the lookup payloads are $\rn{67}$ values of
$2^b - 1$.  Ordinary multiplies are $\chop{67}$; ordinary adds
$\rn{64}$; the $\ln 2$ scaling and two long-path products are
$\rn{64}$; the terminal operation performs architectural RN/RD/RU
rounding.  $x = \pm 1$ and the documented out-of-range behavior are
exact.  Exhaustive traversals of the binary64-embedded space at
$2^{24}$ through $2^{32}$ scale (up to $6.44\times10^{9}$ results)
show zero finite-value discrepancies; NaN-quieting semantics are
modeled and clean on the standing regressions.

\section{Lineage: the Pentium ROM, its errata, and Zen 3}
\label{sec:lineage}

\paragraph{The constants are the 1993 constants.}  The polynomial
coefficients and table entries that reproduce Skylake bit-for-bit are
the physically decoded Pentium constant ROM entries~\cite{shirriff2025},
evaluated through the $\chop{67}$ multiply datapath.  Intel's
formally verified Itanium-era replacement~\cite{harrison2000} never
displaced this kernel in x86 silicon: on random small inputs the two
algorithms usually coincide, but targeted boundary discriminators
select the Pentium path uniquely.

\paragraph{Two single-bit errata in the published decode.}  With the
ROM appendix's printed values, the reconstruction fails grossly in
exactly two table cells; with single-bit corrections it is exact
there like everywhere else.  Both corrections are recoverable from
Skylake behavior alone, and both are confirmed by direct calculation:

\begin{center}
\small
\begin{tabular}{llllr}
\toprule
row & entry & printed significand & corrected & deviation\\
\midrule
186 & $\cos(44/64)$ & \hex{62ec41e9772401864} & \hex{62ec4169772401864} & $+2^{43}$ sig-units ($+2^{-24}$)\\
189 & $\cos(18/64)$ & \hex{7af8853ddbbe9ffd0} & \hex{7af8853ddbbe9efd0} & $+2^{12}$ sig-units ($+2^{-55}$)\\
\bottomrule
\end{tabular}
\end{center}

\noindent Row 186 is checkable on a calculator: the printed decimal
column reads $0.7728350058$, but $\cos(0.6875) = 0.7728349461\ldots$
---wrong from the seventh decimal place, while all fourteen other
trig entries sit within $\pm0.95$ units of the 68th significand bit.
Since the article's decimal column is derived from the hex, a
single-bit slip in the read-out pipeline explains each; the
alternative---that the die truly holds these bits---would have made
the original Pentium's \instr{FCOS} visibly inaccurate
(${\sim}6\times10^{-8}$) on $[0.625, 0.75)$, contradicting
contemporary accuracy studies.  Modern silicon is thus also a
\emph{witness} for the intended master constants.  The same method
independently corrects one published \instr{F2XM1} table payload (the
$b = -57/128$ entry: \hex{43fcb5810d1604f33}; the published
\hex{{\ldots}4f37} variant produces 3{,}869 mode misses on 2{,}563
cell inputs, the corrected value zero).

\paragraph{AMD Zen 3 runs the same family.}  A volunteer-run capture
kit executed the identical input corpora on a Zen~3 core (Ryzen~5
5600H).  Findings: (i) the Skylake behavioral model matches Zen~3
bit-for-bit on 94.3\% of the 240k dense set and 96.2\% of the
structured sweep, with \emph{every} disagreement exactly 1~$\ulp$ on
one output---same algorithm family, slightly different terminal
resolution; (ii) on 8{,}000 operands within a few $\ulp$ of
$k\pi/2$---where a single differing $\pi$ bit would diverge by many
$\ulp$s---Zen~3 is \emph{8{,}000/8{,}000 bit-identical} to Skylake:
the 66-bit-$\pi$ reduction is shared exactly, falsifying (for Zen~3)
the folklore that AMD's x87 trigonometry is more accurate than
Intel's; (iii) on AMD, \instr{FSIN}/\instr{FCOS} and \instr{FSINCOS}
agree on 240{,}000/240{,}000 operands (a single internal path),
whereas Skylake's standalone and paired paths diverge on 987 of
them---a clean microarchitectural fingerprint distinguishing the
vendors.

\section{Validation}
\label{sec:validation}

\subsection{Standalone promotion regression}

H1707 checks the packaged fixed program and H1708 checks its promotion into
the default source. The latter scores 72 authenticated retained bank/mode
inventories against raw hardware outputs and applicable \texttt{C1}:
\begin{center}
\begin{tabular}{lrrr}
\toprule
Path & Output appearances & \texttt{C1} checks & Misses\\
\midrule
Polynomial & 723{,}546 & 723{,}546 & 0\\
Table & 2{,}472{,}906 & 2{,}472{,}906 & 0\\
Tiny & 182{,}535 & 182{,}535 & 0\\
Special/range & 30 & 0 & 0\\
\midrule
Total & 3{,}379{,}017 & 3{,}378{,}987 & 0\\
\bottomrule
\end{tabular}
\end{center}
The 81 external frontier rows over 79 operands all match, as do 53 retained
legacy tuples. These are failures of the preceding model, not 81 remaining
failures of the promoted algorithm. Some legacy/frontier tuples overlap.
Bank appearances also overlap earlier audits and are not unique-observation
counts. The original raw-stream and positional-provenance limitations are
retained in the audit records; historical tuple excerpts do not acquire
full-stream provenance merely by being rechecked.

The polynomial, table and tiny headers are byte-identical in arithmetic to
the validated experimental versions after removing only the added C1 export,
opt-in trace guard and promotion comment. O0, O2, O3 and UBSan builds agree
with the archived candidate on a mixed encoding/domain software bank.
Default diagnostic output is quiet. H1708 preserved the then-current paired
path; the later paired promotions change only that path. Their regressions
rechecks all displayed standalone outputs and applicable C1 indicators,
including all 81 frontier rows, with zero misses.

\subsection{Paired promotion and fresh validation}
\label{sec:paired-validation}

The all-product program repairs three sine-lane misses over two operands in
the archived paired program, with no
regressions across 15 retained banks: 11{,}919{,}273 instruction/RC appearances,
23{,}838{,}534 lane results, six C2 responses and 11{,}919{,}267 applicable
cosine-bound C1 checks. The retained modes are RN/RD/RU. StageA streams have
snapshot/positional provenance rather than original checksum manifests; this
limitation is preserved. These overlapping appearances are not unique tuples.

A 20-million-input software search found 38 materialization discriminators.
Signed counterparts, neighbors, exact reduced preimages and domain/boundary
controls formed 1{,}182 proposed operands. Local history checks excluded 32
operands before freezing the remaining 1{,}150, including 138 discriminators.
The independent integer/rational implementation agreed with every prediction.
No hardware label was used to adjust the candidate after the freeze.

H1712 executed each of the resulting 13{,}800 tuples once on the reference
Xeon: \instr{FSINCOS} at all four RC modes and PC24/53/64. The numerical and
status results are:
\begin{center}
\small
\begin{tabular}{lrrr}
\toprule
Check & Cases & Predecessor misses & Promoted misses\\
\midrule
Sine output & 13{,}728 & 558 & 0\\
Cosine output & 13{,}728 & 150 & 0\\
External cosine C1 & 13{,}728 & 66 & 0\\
C2 response & 13{,}800 & 0 & 0\\
C2 operand preservation & 72 & 0 & 0\\
\bottomrule
\end{tabular}
\end{center}
The predecessor fails 750 distinct tuples; the columns overlap. Every
predeclared discriminating tuple selects the new prediction. The 150 cosine
failures arise when quadrant mapping exposes the changed internal sine arm.
Route counts are 11{,}400 polynomial, 1{,}584 table, 744 tiny and 72 C2 tuples.
Initial operand/control, final control word and expected stack mapping also
pass. Undefined condition bits and unused register payloads are not scored.

The campaign includes 33 exact positive reduced preimages and their signed
counterparts. Its 240 nearby reduction-boundary controls are brackets, not
exact isomorphs. Both the last-product and all-product schedules passed this
challenge; it did not distinguish their earlier Horner edges.

This separator distinguishes the tested schedules and selects all-product
materialization. H1717
promotes that fixed schedule into the default executable. Replaying the
15 retained banks and immutable H1712 outputs gives zero numerical or
applicable C1 misses. Independent rational/integer and O0/O2/O3/UBSan checks
verify the promoted implementation and preservation of standalone behavior.
Fresh evidence belongs to the original observations, not later compiler
builds or replays; finite agreement does not establish all-input equivalence
to silicon.

The promoted executable also matches the i7 review census of
45{,}517{,}233 paired instruction/RC rows (91{,}034{,}466 lanes and
45{,}517{,}233 C1 checks), including the separating operand. This three-bank
census overlaps the retained banks and is not added to them as a unique-input
count. Replays of the saved H1712/H1714/H1715 and independent-review captures
match all 357{,}360 instruction tuples across the three instructions:
482{,}304 output lanes and 353{,}520 applicable C1 checks. No captures are
repeated for promotion. These are regressions on the recorded CPU contexts,
not evidence of parity across every Intel generation.

\subsection{Independent and prospective evidence}

Independent rational implementations recomputed all 578{,}292 polynomial
and 2{,}276{,}661 table output/C1 appearances in the original component
audits. A separate integer tiny formula checked 109{,}905 tiny appearances.
These component counts overlap the promotion regression.

Three previously frozen, one-shot challenges provide prospective evidence:
\begin{center}
\small
\begin{tabular}{p{0.49\linewidth}rr}
\toprule
Challenge & Outputs & Predicted \texttt{C1}\\
\midrule
H1624: direct cosine, four RC modes & 1{,}668 & 1{,}668\\
H1641: boundaries, reductions, encodings, RC/PC & 6{,}432 & 4{,}512\\
H1694: remaining signed exact centers, RC/PC & 624 & 624\\
\bottomrule
\end{tabular}
\end{center}
Every listed frozen prediction passed, without changing the numerical
candidate after observation. H1624 selected 417 eligible operands from a
label-free adversarial search, including model separators and rounding-boundary
controls; it is confined to direct positive cosine in $[1/8,1/4)$.
H1641 covered both instructions and PC24/PC53/PC64 and also passed 3{,}744
predicted exception masks. H1694 combined with 48 already-observed H1641
center tuples covers all 672 keys in the specified finite exact-center
universe: 28 signed external operands, two instructions, four RC modes and
three PC settings. Unknown status fields were not counted as predicted passes.
These banks were not re-executed for promotion.

The table negative control $M_R\mapsto\rn{64}\circ M$ fails 102 retained
outputs and 139 C1 appearances. Hence the destination-rounding distinction
is observable, not algebraic notation chosen after the fact.

Domain and implementation checks supplement these measurements: exact
reduction bounds and modular certificates, conversion/bypass grid arguments,
integer helper checks, accumulator-width bounds and configured-route audits.
They justify the specified C computation under its entry contracts, without
turning a software equivalence theorem into a claim about hidden silicon.

\subsection{Historical evidence is not current validation credit}
\label{sec:residual}

The older R96-stage suite and its captured overlays
remain historical records. That entire suite is \emph{not} claimed here
as a new regression of the promoted source. Likewise, generated-but-uncaptured
software candidates do not count as hardware tests.
R96 remains an incomplete empirical predecessor; a high cache success rate
never made it a general solution.

The rejected-selector record includes QX falsifications on 31/31 challenge
pairs and Q falsifications on 75/75. A 255-representation audit found no
candidate surviving the control and adversarial walls. Complete coarse
12-phase carry histories contain positive/negative collisions, ruling out
deterministic selectors driven solely by those histories. These are bounded
negative results about specified model classes, not proof that the instruction
is non-deterministic or lacks a general arithmetic representation.
The successful standalone and paired programs change the arithmetic graph instead of
adding a selector to those failed models. Detailed exploration and abandoned
fits belong to the research handoff, not the current algorithm.

\section{Discussion}
\label{sec:discussion}

The result is a validated, general numerical reconstruction of standalone
\instr{FSIN}/\instr{FCOS} and paired \instr{FSINCOS} for the selected Skylake
reference, with no known
mismatch in the stated evidence. Its piecewise structure comes from encoding
classes, exact range reduction, native table cells and magnitude thresholds.
The small coefficient ROM is algorithm data, not a lookup of known failures.

The evidence does not establish a unique physical datapath: mathematically
equivalent operand orderings and implementations can yield the same outputs.
It also does not warrant cross-generation parity. Earlier Zen~3 comparisons
show shared algorithmic lineage but last-bit differences.

Complete pointer-register emulation, undefined condition bits and obscure
restore-history behavior are outside the numerical claim; supported numerical
rounding/control behavior is not replaced by a claim of complete x87 state
emulation. Future counterexamples must be retained and used to revise the
general arithmetic, not hidden with an operand-specific patch.

\section{Availability}

The main C implementation is \path{src/fsincos_skylake.c}.
The fixed kernels reside in \path{src/general/}: the three
\texttt{standalone\_*} headers and \texttt{paired.h}.
The standard build enables the validated paths for all three instructions.
A restricted standalone CLI is available through the
\path{general-candidate} build target. The historical translation unit still
contains inactive experimental bodies; they are not part of the promoted
standalone or paired numerical routes. Further source isolation and embedding
API design are optional packaging work, not evidence for a new arithmetic rule.

The H1707/H1708 and H1710--H1718 manifests and reports preserve source, binary, input and
capture hashes and all regression results. The component specifications,
independent verifiers and immutable H1624/H1641/H1694/H1712 campaign records are
retained with them. The pre-promotion manuscript and source are archived;
earlier research notes remain historical and are not silently rewritten
as prospective successes.

\subsection*{Acknowledgments}

This work builds directly on Ken Shirriff's physical decode of the
Pentium FPU constant ROM.  The experimental campaign, analysis, and
drafting were conducted with substantial AI assistance (Anthropic's
Claude and OpenAI Codex); all hardware measurements are from the
author's capture infrastructure and volunteer-contributed
cross-vendor captures.

\clearpage
\appendix
\section{Programmer's pseudocode}
\label{app:pseudocode}

\subsection{Entry, reduction and dispatch}

The following Python-like code describes the numerical values of
\instr{FSIN}, \instr{FCOS} and \instr{FSINCOS}. All ordinary arithmetic is
\emph{exact} signed integer/dyadic arithmetic; only named rounding operations
discard bits. In particular, \texttt{**} means exponentiation and
\texttt{divmod} gives integer quotient and remainder. Host \texttt{double}
arithmetic would change the program. \texttt{x80} denotes an exact x87 value;
\texttt{signbit} returns 0 or 1; \texttt{mod 4} below is
nonnegative even for negative quadrants.

\begin{lstlisting}[style=pseudocode]
def reduce(x):
    # Direct-path boundary as an exact x87 magnitude.
    if abs(x) < x80(0x3ffe, 0xc90fdaa22168c234):
        return x, 0
    s, e = normalized_parts(abs(x))       # abs(x) = s * 2**(e-63)
    M66 = 0x3243f6a8885a308d3
    dividend = s << (e + 2)
    q, remainder = divmod(dividend, M66)
    if 2 * remainder > M66: q += 1        # No halfway tie: M66 is odd.
    residual = (dividend - q * M66) * 2**(-65)
    if signbit(x): return -residual, -q
    return residual, q

def trig(op, x, rc):
    if special_encoding_or_zero(x):
        return special_values(op, x)     # Section 4.1; before reduction.
    if abs(x) >= 2**63: return C2_operand_unchanged(x)
    phases = [0, 1] if op == FSINCOS else [0 if op == FSIN else 1]
    y, c1 = {}, {}
    if abs(x) < 2**(-68):                 # External-input bypass.
        return select_lanes(phases, [x, 1]), 0
    residual, quadrant = reduce(x)
    r = abs(residual)
    if 2**(-32) <= r < 1/4 and op == FSINCOS:
        pair = paired_poly(r)            # Compute once, not two FSIN/FCOS calls.
    for phase in phases:                 # 0: external sine; 1: external cosine.
        n = (quadrant + phase) % 4
        cosine = (n & 1) != 0
        negative = (n >> 1) ^ (0 if cosine else signbit(residual))
        if r < 2**(-32):
            y[phase], c1[phase] = tiny(r, cosine, negative, rc)
        else:
            if r >= 1/4: u = table(r, cosine)
            elif op == FSINCOS: u = pair[int(cosine)]
            else: u = standalone_poly(r, cosine)
            y[phase], c1[phase] = finish(u, negative, rc)
    return y, c1[phases[-1]]              # FSINCOS C1 comes from external cosine.
\end{lstlisting}

\noindent Here \texttt{y} and \texttt{c1} are maps keyed by the requested
phase. The returned lanes are mathematical sine/cosine values, not an x87
stack or exception-state API. Special-value handling follows
Section~\ref{sec:specials}; its undefined condition bits are not specified
by this pseudocode.

\clearpage
\subsection{Standalone and paired polynomial kernels}

\texttt{T67} and \texttt{T64} truncate the magnitude to 67 and 64 significant
bits, preserving sign. \texttt{RN64} rounds once to nearest-even at 64 bits.
\texttt{RC64} uses the requested architectural rounding mode, also at 64 bits:
nearest-even (RN), down (RD), up (RU), or toward zero (RZ).
The following operators retain their operand order:

\begin{lstlisting}[style=pseudocode]
def M(x, y):  return T67(T67(x) * T64(y))
def A(x, y):  return RN64(x + y)
def MR(x, y): return RN64(T67(x) * T64(y))  # One rounding of exact port product.
\end{lstlisting}

\noindent \texttt{MR(x,y)} cannot be replaced by \texttt{RN64(M(x,y))}.
Coefficients use one-based indices: \texttt{S6}, \texttt{C6}, \texttt{S4}
and \texttt{C4} are the signed native ROM families in
\path{src/p5_rom_constants.h}, not Taylor-series approximations.

\begin{lstlisting}[style=pseudocode]
def standalone_poly(r, cosine):
    K = C6 if cosine else S6
    square = M(r, r)
    fourth = M(square, square)            # Includes T64 on the second port.
    odd = M(fourth, K[5])
    odd = A(K[3], odd)
    odd = A(K[1], M(fourth, odd))
    even = M(fourth, K[6])
    even = A(K[4], even)
    even = A(K[2], M(fourth, even))
    left = M(square, odd)
    right = M(fourth, even)
    if cosine: return 1 + T67(left + right)
    return r + M(r, A(left, right))

def paired_poly(r):
    square = T67(r * r)
    p, q = S6[6], C6[6]
    for i in [5, 4, 3, 2, 1]:
        p = RN64(T67(p * square) + S6[i]) # Every product is cut before RN64.
        q = RN64(T67(q * square) + C6[i]) # Including the last on BOTH arms.
    sine = r + T67(RN64(p * square) * r)
    cosine = 1 + T67(q * square)
    return [sine, cosine]
\end{lstlisting}

\noindent Each kernel returns an unrounded positive-residual prevalue.
The caller selects the internal sine or cosine branch and applies its
quadrant sign \emph{before} final architectural rounding. Earlier internal
rounding stays fixed regardless of RC. The standalone split into odd and
even coefficient indices and the paired Horner loop are different numerical
schedules; substituting one for the other changes observable outputs.

\clearpage
\subsection{Shared table kernel, tiny results and final rounding}

The table is indexed by center numerator \texttt{b}. \texttt{ROM[b]} contains
the corrected native sine/cosine words, not fresh host-library evaluations.
Cell selection and center subtraction remain exact at binary boundaries.

\begin{lstlisting}[style=pseudocode]
def horner4(square, K):
    v = K[4]
    for i in [3, 2, 1]: v = A(K[i], M(square, v))
    return v

def table(r, cosine):
    if r < 1/2: b = 18 + 4 * floor(16 * (r - 1/4))
    else: b = 36 + 8 * min(2, floor(8 * (r - 1/2)))
    a = r - b/64
    square = M(a, a)
    K = copy(S4)
    K[4] -= 2**(-25)                      # Validated coefficient adjustment.
    p = horner4(square, K)
    q = horner4(square, C4)
    v = A(a, M(M(square, p), a))
    w = MR(square, q)
    ts, tc = ROM[b]
    if cosine: return tc + T67(-M(ts, v) + M(tc, w))
    return ts + T67(M(tc, v) + M(ts, w))

def tiny(r, cosine, negative, rc):
    lead = 1 if cosine else r
    toward_zero = (rc == RZ or (rc == RD and not negative)
                   or (rc == RU and negative))
    if toward_zero: lead = pred64(lead)
    return (-lead if negative else lead), int(not toward_zero)

def finish(u, negative, rc):
    signed = -u if negative else u
    y = RC64(signed, rc)
    return y, int(abs(y) > abs(signed))   # C1 is a magnitude increment.
\end{lstlisting}

\noindent \texttt{pred64(v)} is the greatest nonnegative 64-significand-bit
value strictly less than positive \texttt{v}; at a power of two it crosses
into the preceding binade. It is used only for non-bypass tiny values.
RN and directed rounding away from zero keep the leading magnitude and set
\texttt{C1}; the earlier external-input bypass instead returns \texttt{C1=0}
in every mode. \texttt{finish} measures only final rounding, not whether any
earlier operation was inexact.

The kernel headers reside in \path{src/general/}.
\texttt{standalone\_polynomial.h} and \texttt{paired.h} supply the polynomial
paths; \texttt{standalone\_table.h} and \texttt{standalone\_tiny.h} supply
the shared paths. The exact reducer and encoding helpers are in
\path{src/fsincos_skylake.c}.
This walkthrough restates the validated numerical program and its evidence
limits; it does not extend the claim to full x87 state emulation.

\clearpage
\begin{thebibliography}{9}
\scriptsize
\raggedright
\itemsep0pt

\bibitem{shirriff2025}
K.~Shirriff.
\newblock Pi in the Pentium: reverse-engineering the constants in its
  floating-point unit.
\newblock \url{https://www.righto.com/2025/01/pentium-floating-point-ROM.html},
  January 2025.

\bibitem{shirriff-microcode2025}
K.~Shirriff.
\newblock Notes on the Pentium's microcode circuitry.
\newblock \url{https://www.righto.com/2025/03/pentium-microcde-rom-circuitry.html},
  March 2025.

\bibitem{bosch-p6-1}
P.~Bosch.
\newblock Pentium II microcode: collected tweets, part 1.
\newblock \url{https://pbx.sh/pentiumii-part1/}, November 2022.

\bibitem{bosch-p6-2}
P.~Bosch.
\newblock Pentium II microcode: collected tweets, part 2.
\newblock \url{https://pbx.sh/pentiumii-part2/}, November 2022.

\bibitem{bosch-p6tools}
P.~Bosch.
\newblock \texttt{p6tools}: Pentium II microcode (dis)assembler and
  (de)scrambler.
\newblock \url{https://github.com/peterbjornx/p6tools}.

\bibitem{dawson2014}
B.~Dawson.
\newblock Intel underestimates error bounds by 1.3 quintillion.
\newblock \url{https://randomascii.wordpress.com/2014/10/09/intel-underestimates-error-bounds-by-1-3-quintillion/},
  October 2014.

\bibitem{intel-sdm}
Intel Corporation.
\newblock \emph{Intel 64 and IA-32 Architectures Software Developer's
  Manual}, Vol.~1, \S8.3.8 (transcendental instruction accuracy).

\bibitem{morrow1997}
W.~R. Morrow.
\newblock Method and apparatus for rounding operands using previous
  rounding history.
\newblock U.S. Patent 5,612,909, assigned to Intel Corporation, March 1997.

\bibitem{palaniswami1993}
K.~J. Palaniswami.
\newblock Computation of sticky-bit in parallel with partial products in a
  floating point multiplier unit.
\newblock U.S. Patent 5,260,889, assigned to Intel Corporation, November 1993.

\bibitem{harrison2000}
J.~Harrison.
\newblock Formal verification of floating point trigonometric functions.
\newblock In \emph{Formal Methods in Computer-Aided Design (FMCAD)},
  LNCS 1954, pp. 217--233, 2000.

\bibitem{harrison1999}
J.~Harrison, T.~Kubaska, S.~Story, and P.~T.~P. Tang.
\newblock The computation of transcendental functions on the IA-64
  architecture.
\newblock \emph{Intel Technology Journal}, Q4 1999.

\bibitem{lynch1995}
T.~Lynch, A.~Ahmed, M.~Schulte, T.~Callaway, and R.~Tisdale.
\newblock The K5 transcendental functions.
\newblock In \emph{Proc. 12th IEEE Symposium on Computer Arithmetic
  (ARITH-12)}, pp. 163--170, 1995.

\bibitem{nicely1994}
T.~R. Nicely.
\newblock Pentium FDIV flaw, 1994.
\newblock Public disclosure e-mail and subsequent analyses.

\bibitem{coe1995}
T.~Coe and P.~T.~P. Tang.
\newblock It takes six ones to reach a flaw.
\newblock In \emph{Proc. 12th IEEE Symposium on Computer Arithmetic
  (ARITH-12)}, pp. 140--146, 1995.

\bibitem{muller2016}
J.-M. Muller.
\newblock \emph{Elementary Functions: Algorithms and Implementation}.
\newblock Birkh\"auser, 3rd edition, 2016.

\end{thebibliography}

\end{document}
